% ============================================================================
\section{Proofs and verification}
\label{app:proofs}

This appendix proves that the closed forms \eqref{eq:potentials} solve the Helmholtz--Hodge decomposition, proves the
two other propositions and gives the closed forms that the code evaluates.

\begin{proposition}[Closed-form Helmholtz--Hodge potentials]\label{prop:hhd}
For a single agent, $\nabla\big(u\,\kappa_\sigma(\vd)\big)-\nabla^{\perp}\big(w\,\kappa_\sigma(\vd)\big)
=\vv\,e^{-\|\vd\|^2/2\sigma^2}$ at every $\vd\neq\vec 0$. Consequently the potentials \eqref{eq:potentials} satisfy,
at every point $\vr$,
\[
\nabla\varphi_\sigma+\nabla^{\perp}\psi_\sigma=\frac{1}{NT\sigma\vbar}\sum_{t}\vV_\sigma(\vr,t).
\]
\end{proposition}

\begin{proof}[Proof of \cref{prop:hhd}]
The idea is to invert the Laplacian on the flux of one agent, that is, to solve the Poisson equation of the plane
for a vector field. The solution is then split into a gradient and a rotated gradient, and the two potentials are
the scalar fields that appear in this split.

Consider one agent at the origin, with a constant velocity $\vv$, and write $\varrho=\|\vd\|$ for the
distance to it. Its contribution to the flux is $\vv\,e^{-\varrho^2/2\sigma^2}$, a Gaussian bump carried in
the direction of $\vv$. We first look for a radial function $\Gamma(\varrho)$ whose Laplacian is this
Gaussian, up to a factor $1/\sigma^2$. For a radial function, the Laplacian reads
$\Delta\Gamma=\varrho^{-1}(\varrho\Gamma')'$. Integrating from the origin, where $\Gamma$ must stay smooth, gives
\[
\varrho\,\Gamma'(\varrho)=\int_0^{\varrho}\frac{s}{\sigma^2}\,e^{-s^2/2\sigma^2}\,\mathrm{d}s=1-e^{-\varrho^2/2\sigma^2},
\]
which gives
\[
\Delta\Gamma=\frac{1}{\varrho}\,\frac{\mathrm{d}}{\mathrm{d}\varrho}\Big(1-e^{-\varrho^2/2\sigma^2}\Big)
=\frac{1}{\sigma^2}\,e^{-\varrho^2/2\sigma^2}.
\]
Far from the agent, $\Gamma'\simeq1/\varrho$, so $\Gamma$ grows like the logarithm of the distance, as the free-space
Green's function of the plane does \citep{bhatia2014natural}. Near the agent, the Gaussian smooths it.
Compared with \eqref{eq:kappa}, this choice reads $\sigma^2\Gamma'(\varrho)=\varrho\,\kappa_\sigma(\varrho)$:
the kernel $\kappa_\sigma$ is $\sigma^2$ times the radial derivative of $\Gamma$, divided by the
distance.

Second, every smooth vector field $\vW=(W_x,W_y)$ of the plane satisfies
\[
\Delta\vW=\nabla(\nabla\!\cdot\vW)+\nabla^{\perp}(\nabla^{\perp}\!\cdot\vW),\qquad
\nabla^{\perp}\!\cdot\vW=\partial_xW_y-\partial_yW_x,
\]
as the expansion of both sides component by component shows. We apply this identity to
$\vW=\sigma^2\Gamma\,\vv$. Since $\vv$ is constant, $\Delta\vW=\sigma^2(\Delta\Gamma)\,\vv=\vv\,e^{-\varrho^2/2\sigma^2}$,
which is the contribution of the agent. The identity splits it into the gradient of
$\Phi=\nabla\!\cdot\vW=\sigma^2\,\vv\cdot\nabla\Gamma$ and the rotated gradient of
$\Psi=\nabla^{\perp}\!\cdot\vW=\sigma^2\big(v_y\,\partial_x\Gamma-v_x\,\partial_y\Gamma\big)$:
\[
\vv\,e^{-\varrho^2/2\sigma^2}=\nabla\Phi+\nabla^{\perp}\Psi .
\]

Third, we write $\Phi$ and $\Psi$ with the kernel. The gradient of a radial function points along $\vd$,
and $\sigma^2\nabla\Gamma=\sigma^2\Gamma'(\varrho)\,\vd/\varrho=\kappa_\sigma(\vd)\,\vd$. Hence
$\Phi=\kappa_\sigma(\vd)\,(\vv\cdot\vd)=u\,\kappa_\sigma(\vd)$ and
$\Psi=\kappa_\sigma(\vd)\,(v_yd_x-v_xd_y)=-w\,\kappa_\sigma(\vd)$, which is the single-agent identity of the
proposition. The computation holds wherever $\vd\neq\vec 0$. At the agent itself both sides are smooth, and
the identity follows by continuity.

Finally, the identity is linear in the velocity. Summing it over the agents and the instants, with the
factor $1/(NT\sigma\vbar)$, gives the potentials \eqref{eq:potentials} on the left and
$\frac{1}{NT\sigma\vbar}\sum_t\vV_\sigma(\vr,t)$ on the right, since $\vV_\sigma=\sum_i\vv_iO_i$. The two
potentials therefore rebuild the time-averaged flux exactly, without a grid and without a boundary
condition.
\end{proof}

\begin{proposition}[Invariance]\label{prop:so2}
Let a rigid transformation of the plane, a rotation $Q\in\mathrm{SO}(2)$ followed by a translation, act
on all positions, on all velocities and on the evaluation point. Then $\rho_\sigma$, $K_\sigma$, $\varphi_\sigma$,
$\psi_\sigma$ and $J_\sigma$ are unchanged, while $\nabla\varphi_\sigma$ and $\nabla\psi_\sigma$
are multiplied by $Q$.
\end{proposition}
\begin{proof}
Every field is built from a few elementary quantities: the distances $\|\vd_i\|$ between the point and
the agents, the speeds $\|\vv_i\|$, the scalar products $u_i=\vv_i\cdot\vd_i$, the determinants
$w_i=\det[\vv_i,\vd_i]$, and the averages $\vbar$ and $\overline{v^2}$. A translation changes neither $\vd_i$
nor $\vv_i$. A common rotation preserves lengths, angles and the orientation of the plane, so it changes
none of these quantities either. Hence $\rho_\sigma$, $K_\sigma$, $\varphi_\sigma$ and $\psi_\sigma$ take
the same value at the transformed point, and their gradients turn with the scene: they are multiplied
by $Q$. The coupling keeps its value, since a rotation does not change a determinant:
$\det[Q\nabla\varphi_\sigma,Q\nabla\psi_\sigma]=\det Q\,\det[\nabla\varphi_\sigma,\nabla\psi_\sigma]
=\det[\nabla\varphi_\sigma,\nabla\psi_\sigma]$, because $\det Q=1$.
\end{proof}
The seeds are selected from the agents and from the fields themselves, so the diagrams inherit this
invariance, up to floating-point effects in the choice of the seeds, which \cref{app:invariance}
measures.

\begin{proposition}[Tail of the footprint]\label{prop:tail}
For $O_i(\vr)=\exp(-\|\vr-\vx_i\|^2/2\sigma^2)$ and any $\varrho\ge0$,
$\int_{\|\vr-\vx_i\|>\varrho}O_i(\vr)\,\mathrm{d}\vr=2\pi\sigma^2e^{-\varrho^2/2\sigma^2}$. At $\varrho=5\sigma$
the fraction of the mass outside the evaluation window is $e^{-12.5}\approx3.7\times10^{-6}$.
\end{proposition}
\begin{proof}
The footprint depends only on the distance $s=\|\vr-\vx_i\|$ to the agent, so we integrate in polar
coordinates centered on it, where the element of area is $2\pi s\,\mathrm{d}s$:
\[
\int_{\|\vr-\vx_i\|>\varrho}O_i(\vr)\,\mathrm{d}\vr=\int_\varrho^\infty e^{-s^2/2\sigma^2}\,2\pi s\,\mathrm{d}s
=2\pi\sigma^2\Big[-e^{-s^2/2\sigma^2}\Big]_{\varrho}^{\infty}=2\pi\sigma^2e^{-\varrho^2/2\sigma^2}.
\]
With $\varrho=0$ the whole mass is $2\pi\sigma^2$, so the fraction left outside a window of radius
$\varrho$ is $e^{-\varrho^2/2\sigma^2}$. At $\varrho=5\sigma$ it is $e^{-25/2}=e^{-12.5}\approx3.7\times10^{-6}$.
\end{proof}

The coupling field needs the gradients of the two potentials. They follow from the same closed forms,
by differentiating $u_i\kappa_\sigma(\vd_i)$ and $w_i\kappa_\sigma(\vd_i)$ with respect to the point. With
$c=1/(NT\sigma\vbar)$, $\vv_i^{\perp}=(-v_{i,y},v_{i,x})$ and
$\lambda_i=\big(e^{-\|\vd_i\|^2/2\sigma^2}-2\kappa_\sigma(\vd_i)\big)/\|\vd_i\|^2$,
\[
\nabla\varphi_\sigma=c\sum_{t,i}\big(\kappa_\sigma(\vd_i)\,\vv_i+\lambda_i\,u_i\,\vd_i\big),\qquad
\nabla\psi_\sigma=-c\sum_{t,i}\big(\kappa_\sigma(\vd_i)\,\vv_i^{\perp}+\lambda_i\,w_i\,\vd_i\big).
\]
A seed taken from the agent positions coincides with an agent at some instant, so these expressions must remain
accurate at $\vd_i=\vec 0$. Two details ensure it. The code computes $1-e^{-s}$ as $-\operatorname{expm1}(-s)$,
which keeps the relative precision of $\kappa_\sigma$ as $\|\vd\|\to0$. It also floors $\|\vd\|^2$ at
$10^{-12}$\,m$^2$, where $\kappa_\sigma$ equals $1/2$ to better than the single precision in which the
fields are stored. The factor $\lambda_i$ tends to $-1/(4\sigma^2)$ as $\vd_i\to\vec 0$, but its numerator is
the difference of two numbers close to one, which loses precision. It is multiplied by $u_i\,\vd_i$ or
$w_i\,\vd_i$, of order $\|\vd_i\|^2$, so this loss does not reach the gradients. The derivatives with respect
to the positions (\cref{sec:res_attribution}) are taken by automatic differentiation of the same
expressions.

Every statement of this appendix is also checked by computer. The script
\nolinkurl{manuscript/verify_propositions.py} verifies with SymPy the single-agent identity of
\cref{prop:hhd}, the orthogonality relations $\nabla\cdot\nabla^{\perp}\psi=0$ and
$\nabla^{\perp}\cdot\nabla\varphi=0$, the invariance of the elementary terms of \cref{prop:so2}, and the
tail of \cref{prop:tail}, which a numerical quadrature confirms to $\nm{tail-quad-err}$. It also checks
the gradients above, the limit of $\lambda_i$ and their value at an agent, and the rotation invariance of
the magnitude of the traceless part of $\tM_2$ (\cref{sec:representation}). On a synthetic milling group, it
checks numerically the envelope property of \cref{app:attribution}.
